\documentclass[12pt,a4paper]{report}
\usepackage[utf8]{inputenc}
\usepackage{amsmath}
\usepackage{amsfonts}
\usepackage{amssymb}
\usepackage{amsthm}
\usepackage{hyperref}

\usepackage{multicol}
\usepackage{fancyhdr}
\usepackage[inline]{enumitem}
\usepackage{tikz}
\usepackage{tikz-cd}
\usetikzlibrary{calc}
\usetikzlibrary{shapes.geometric}
\usepackage[margin=0.5in]{geometry}
\usepackage{xcolor}
\usepackage{ esint }

\hypersetup{
    colorlinks=true,
    linkcolor=blue,
    filecolor=magenta,      
    urlcolor=cyan,
    pdftitle={PDF2 Sample Test},
    pdfpagemode=FullScreen,
    }

%\urlstyle{same}

\newcommand{\CLASSNAME}{Comprehensive Exam -- PDE2}
\newcommand{\STUDENTNAME}{Paul Carmody}
\newcommand{\ASSIGNMENT}{NOTES }
\newcommand{\DUEDATE}{April 29, 2026}
\newcommand{\SEMESTER}{Spring 2026}
\newcommand{\SCHEDULE}{3:00 -- 5:00}
\newcommand{\ROOM}{Crawford 332}

\newcommand{\MMN}{M_{m\times n}}
\newcommand{\FF}{\mathcal{F}}

\pagestyle{fancy}
\fancyhf{}
\chead{ \fancyplain{}{\CLASSNAME} }
%\chead{ \fancyplain{}{\STUDENTNAME} }
\rhead{\thepage}
\fancyfoot[C]{-- \thepage --}
\input{../../MyMacros}

\newcommand{\RED}[1]{\textcolor{red}{#1}}
\newcommand{\BLUE}[1]{\textcolor{blue}{#1}}
\newcommand{\FFF}{\mathcal{F}}

\begin{document}

\begin{center}
	\Large{\CLASSNAME -- \SEMESTER} \\
	\large{ w/Faculty}
\end{center}
\begin{center}
	\STUDENTNAME \\
	\ASSIGNMENT -- \DUEDATE\\
\end{center} 

\begin{enumerate}

\item Let $U = \{x\in \R^n, |x| > 1\}$ and let $u(x) = 1/|x|^\alpha$ for $x \in U$, where $\alpha > 0$.  For which values of $p$ is $U \in W^{1,p}(U)$.?

\HLINE
\textbf{Notes}
\begin{description}
\item $W^{k,p}$ is a set of functions, $u\in L^p$ and $Du\in L^p$.
\item Be conscious fo $|x|$ realizing that polar coordinates are often better.
\item Remember convergences theorems.
\end{description}

%1%%%%%%%%%%%%%%%%%%%%
\HLINE
\begin{itemize}
	\item What values of $n$ and $p$ make $u \in L^p$?
\begin{align*}
	\int_U |u(x)|^p\,dx &< \infty \\
		&= \int_U|x|^{-\alpha p} \, dx\\
		&= C \int_0^\infty r^{n-1-\alpha p}\,dr\quad\text{polar coordinates}
\end{align*}which converges iff $n-1-\alpha p < -1$ or $\alpha p > n$.
	\item What values of $n$ and $p$ make $Du \in L^p$?
	\begin{align*}
		u(x) &= |x|^{-\alpha} \\
		Du(x) &=  -\alpha|x|^{-\alpha-2}x =  -\alpha|x|^{-\alpha-1}  \\
		|Du(x)|^p &=  -\alpha^p|x|^{-p(\alpha+1)}  \\
		\int_U |Du(x)|^p\,dx &= C\int_1^\infty r^{n-1-p(\alpha+1)}\,dr,\quad\text{polar coordinates}
	\end{align*}which converges when $p(\alpha+1)>n$. Therefore, whenv $p> \frac{n}{\alpha}$.
\end{itemize}

%2%%%%%%%%%%%%%%%%%%%%
\newpage
\item Let $U$ be a bounded open subset of $\R^n$ with $\partial U \in C^1$.  Let $c \in L^\infty(U)$ and let $f \in L^2(U)$.  We say that $U\in H^1(U)$ is a weak solution of 
\begin{align*}
	\left \{\begin{array}{rl}
		 -\Delta u+c(x)=f & \text{in } U\\
		\PART{u}{\nu} = 0 & \text{on } \partial U,
	\end{array} \right .
\end{align*}
where $\nu$ is the outward unit normal to $\partial U$, if 
\begin{align*}
	\int_U (\nabla u\cdot\nabla u+cuv)\,dx &= \int_U fv\,dx, \quad \forall v \in H^1(U).
\end{align*}
Assume that 
\begin{align*}
	\inf_{x\in U} c(x) >0.
\end{align*}
Show that there exists a unique weak solution.

\HLINE

To show that the Neumann problem

\begin{align*}
\left\{
\begin{aligned}
-\Delta u + c(x)u &= f &&\text{in } U,\\
\frac{\partial u}{\partial \nu} &=0 &&\text{on } \partial U,
\end{aligned}
\right.
\end{align*}

has a \DEFINE{unique weak solution} $u\in H^1(U$), the standard approach is to apply the \DEFINE{Lax–Milgram Theorem}.

---
\begin{description}

\item \textbf{Step 1: Weak formulation}

Define the bilinear form $B:H^1(U)\times H^1(U)\to \mathbb R$ by

\begin{align*}
B[u,v] = \int_U (\nabla u\cdot \nabla v + c uv)\,dx,
\end{align*}

and define the linear functional $F:H^1(U)\to\mathbb R$ by

\begin{align*}
F(v)=\int_U fv\,dx.
\end{align*}

Then the weak problem is:

\> Find $u\in H^1(U$) such that
\begin{align*}
B[u,v]=F(v)\qquad \forall v\in H^1(U).
\end{align*}

By the \DEFINE{Lax–Milgram theorem}, this has a unique solution provided:

\begin{enumerate}
\item $B$ is bounded,
\item $B$ is coercive,
\item $F$ is bounded.
\end{enumerate}

---

\item \textbf{Step 2: Boundedness of $B$}

We estimate:

\begin{align*}
|B[u,v]| \le \int_U |\nabla u||\nabla v|\,dx+\int_U |cuv|\,dx.
\end{align*}

Using Cauchy–Schwarz,

\begin{align*}
\int_U |\nabla u||\nabla v| \le |\nabla u|_{L^2}|\nabla v|_{L^2}.
\end{align*}

Since $c\in L^\infty(U)$,

\begin{align*}
\int_U |cuv| \le |c|_{L^\infty}|u|_{L^2}|v|_{L^2}.
\end{align*}

Thus,

\begin{align*}
|B[u,v]| \le |\nabla u|_{L^2}|\nabla v|_{L^2} + |c|_{L^\infty}|u|_{L^2}|v|_{L^2}.
\end{align*}

Hence

\begin{align*}
|B[u,v]| \le C|u|_{H^1(U)}|v|_{H^1(U)},
\end{align*}

so $B$ is bounded.

---

\item \textbf{Step 3: Coercivity of $B$}

We compute

\begin{align*}
B[u,u] = \int_U (|\nabla u|^2+c u^2)\,dx.
\end{align*}

Since

\begin{align*}
c(x)\ge c_0>0
\end{align*}

for some constant $c_0$, we get

\begin{align*}
B[u,u] \ge \int_U |\nabla u|^2\,dx+c_0\int_U u^2\,dx.
\end{align*}

Thus,

\begin{align*}
B[u,u] \ge \min(1,c_0)|u|_{H^1(U)}^2.
\end{align*}

Therefore $B$ is coercive.

---

\item \textbf{Step 4: Boundedness of $F$}

We have

\begin{align*}
|F(v)| = \left|\int_U fv\,dx\right| \le |f|_{L^2}|v|_{L^2} \le |f|_{L^2}|v|_{H^1(U)}.
\end{align*}

Thus $F$ is bounded on $H^1(U)$.

---

\item \textbf{Step 5: Apply Lax–Milgram}

Since:

\begin{itemize}
\item $B$ is bounded,
\item $B$ is coercive,
\item $F$ is bounded,
\end{itemize}

the \DEFINE{Lax–Milgram theorem} gives the existence of a unique $u\in H^1(U)$ such that

\begin{align*}
B[u,v]=F(v) \qquad\forall v\in H^1(U).
\end{align*}

That is,

\begin{align*}
\int_U (\nabla u\cdot \nabla v + c uv)\,dx = \int_U fv\,dx \qquad\forall v\in H^1(U).
\end{align*}

Hence there exists a \DEFINE{unique weak solution}.

---

\item \textbf{Final Conclusion}

\begin{align*}
\boxed{\text{There exists a unique weak solution }u\in H^1(U).}
\end{align*}
because the bilinear form
\begin{align*}
B[u,v]=\int_U (\nabla u\cdot\nabla v+cuv)\,dx
\end{align*}
is bounded and coercive on $H^1(U)$, and the linear functional
\begin{align*}
F(v)=\int_U fv\,dx
\end{align*}
is bounded, so the result follows from the Lax–Milgram theorem.

---
\end{description}

If you'd like, I can explain \DEFINE{why the positivity assumption $c(x)\ge c_0>0$ is crucial for uniqueness in the Neumann problem}, since without it uniqueness generally fails.



%3%%%%%%%%%%%%%%%%%%%%
\newpage
\item Let $U$ be a bounded open subset of $\R^n$ with $\partial U \in C^1$.  Let $(f_j) \subset L^2(U)$ converge weakly to $f$ in $L^2(U)$ and let $u_j \in  H_0^1(U)$ be a weak solution of
\begin{align*}
	\left \{\begin{array}{rl}
		 -\Delta u_j=f_j & \text{in } U\\
		u_j = 0 & \text{on } \partial U,
	\end{array} \right .
\end{align*}
Assume that $(u_j)$ is a bounded sequence in $H_0^1(U)$.  Show that a subsequence of $(u_j)$ converges weakly to a weak solution of
\begin{align*}
	\left \{\begin{array}{rl}
		 -\Delta u=f & \text{in } U\\
		u = 0 & \text{on } \partial U,
	\end{array} \right .
\end{align*}

\HLINE

Let $u_j\in H_0^1(U)$ satisfy

\begin{align*}
\int_U \nabla u_j\cdot \nabla v\,dx=\int_U f_jv\,dx
\qquad\forall v\in H_0^1(U),
\end{align*}

and assume:

\begin{enumerate}
\item  $f_j\rightharpoonup f$ weakly in $L^2(U)$,
\item  $(u_j$) is bounded in $H_0^1(U)$.
\end{enumerate}

We want to show:

\> There exists a subsequence $u_{j_k}$ converging weakly in $H_0^1(U)$ to some $u$, and this $u$ is a weak solution of
\begin{align*}
 -\Delta u=f,\qquad u=0\text{ on }\partial U.
\end{align*}

---
\begin{description}

\item \textbf{Step 1: Extract a weakly convergent subsequence}

Since $(u_j)$ is bounded in $H_0^1(U)$, and $H_0^1(U)$ is a Hilbert space, bounded sets are weakly compact.

Therefore, there exists a subsequence $(u_{j_k})$ and some $u\in H_0^1(U)$ such that
\begin{align*}
u_{j_k}\rightharpoonup u
\quad\text{weakly in }H_0^1(U).
\end{align*}

This means in particular,
\begin{align*}
\nabla u_{j_k}\rightharpoonup \nabla u
\quad\text{weakly in }L^2(U).
\end{align*}

---

\item \textbf{Step 2: Pass to the limit in the weak formulation}

For every $v\in H_0^1(U)$,
\begin{align*}
\int_U \nabla u_{j_k}\cdot \nabla v\,dx=\int_U f_{j_k}v\,dx.
\end{align*}

We pass to the limit on both sides.

---

\begin{description}
\item \textbf{Left-hand side}

Because
\begin{align*}
\nabla u_{j_k}\rightharpoonup \nabla u
\quad\text{in }L^2(U),
\end{align*}
we have
\begin{align*}
\int_U \nabla u_{j_k}\cdot \nabla v\,dx
\to
\int_U \nabla u\cdot \nabla v\,dx.
\end{align*}

---

\item \textbf{Right-hand side}

Since $f_{j_k}\rightharpoonup f$ weakly in $L^2(U)$, and $v\in H_0^1(U)\subset L^2(U)$, we get
\begin{align*}
\int_U f_{j_k}v\,dx
\to
\int_U fv\,dx.
\end{align*}

---
\end{description}

\item \textbf{Step 3: Obtain the limiting weak formulation}

Taking limits gives:
\begin{align*}
\int_U \nabla u\cdot \nabla v\,dx = \int_U fv\,dx \qquad\forall v\in H_0^1(U).
\end{align*}

This is exactly the weak formulation of
\begin{align*}
\left\{
\begin{aligned}
-\Delta u&=f &&\text{in }U,\\
u&=0 &&\text{on }\partial U.
\end{aligned}
\right.
\end{align*}

Thus $u$ is a weak solution.

---

\item \textbf{Final conclusion}

There exists a subsequence $u_{j_k}$ such that
\begin{align*}
u_{j_k}\rightharpoonup u
\quad\text{weakly in }H_0^1(U),
\end{align*}
and the limit $u$ satisfies
\begin{align*}
\int_U \nabla u\cdot\nabla v\,dx=\int_U fv\,dx
\qquad\forall v\in H_0^1(U),
\end{align*}
so $u$ is a weak solution of
\begin{align*}
-\Delta u=f,\qquad u=0\text{ on }\partial U.
\end{align*}

---
\end{description}

\begin{align*}
\boxed{\text{A weakly convergent subsequence of }(u_j)\text{ converges to a weak solution }u.}
\end{align*}

This is a standard \DEFINE{weak compactness + passage to the limit} argument in Partial Differential Equations.

If you'd like, I can show why the assumption that $(u_j)$ is bounded in $H_0^1(U)$ is actually automatic from the boundedness of $(f_j$) in $L^2(U)$.

\HLINE

\textbf{Questions}
\begin{itemize}
	\item \textit{Explain `weak convergence'}
	\item \textit{Explain `weak differeentiation'}
	\item \textit{Explain the reltionship between $H_0^1$ and $L^p$ spaces}
\end{itemize}


%4%%%%%%%%%%%%%%%%%%%%
\newpage
\item Let $U$ be a bounded open subset of $\R^n$, $n>3$ with $\partial U \in C^1$.  Let $f: \R\to\R$ be a continous function satisfying
\begin{align*}
	|f(t)| \le c|t|^{q-1}\quad\forall \in \R
\end{align*}
for some constants $c > 0$ and $ 1 \le q < 2$.  We say that $u \in H_o^1(U)$ is a weak solution of 
\begin{align*}
	\left \{\begin{array}{rl}
		 -\Delta u=f(u) & \text{in } U\\
		u = 0 & \text{on } \partial U,
	\end{array} \right .
\end{align*}
if
\begin{align*}
	\int_U \nabla u\cdot\nabla v\, dx = \int_U f(u)v\,dx\quad \forall v \in H_0^1.
\end{align*}
Show that there exists a constant $C >0$ such that 
\begin{align*}
	\|u\|_{H_0^1(U)} \le C
\end{align*}
for all weak solutions $u$.

\HLINE

We want to prove that every weak solution $u\in H_0^1(U)$ of
\begin{align*}
\begin{cases}
-\Delta u=f(u) &\text{in }U,\\
u=0 &\text{on }\partial U,
\end{cases}
\end{align*}
satisfies a \DEFINE{uniform a priori bound}
\begin{align*}
\|u\|_{H_0^1(U)}\le C.
\end{align*}

The key ingredients are:
\begin{enumerate}
\item test the weak equation with $v=u$,
\item use the growth condition on $f$,
\item apply the Sobolev embedding,
\item use the fact that $q<2$.
\end{enumerate}

---
\begin{description}

\item \textbf{Step 1: Use the weak formulation with $v=u$}

Since $u$ is a weak solution,
\begin{align*}
\int_U \nabla u\cdot \nabla v\,dx = \int_U f(u)v\,dx \qquad \forall v\in H_0^1(U).
\end{align*}
Take (v=u). Then
\begin{align*}
\int_U |\nabla u|^2\,dx = \int_U f(u)u\,dx.
\end{align*}
So
\begin{align*}
\|\nabla u\|_{L^2(U)}^2 = \int_U f(u)u\,dx.
\end{align*}

---

\item \textbf{Step 2: Use the growth condition on $f$}

We are given
\begin{align*}
|f(t)|\le c|t|^{q-1}.
\end{align*}
Thus
\begin{align*}
|f(u)u| \le c|u|^q.
\end{align*}
Therefore,
\begin{align*}
\|\nabla u\|_{L^2(U)}^2 \le c\int_U |u|^q\,dx = c\|u\|_{L^q(U)}^q.
\end{align*}
Hence,
\begin{align*}
\|\nabla u\|_{L^2}^2 \le c|\|u\|_{L^q}^{q}.
\end{align*}

---

\item \textbf{Step 3: Use Sobolev embedding}

Since $U$ is bounded and $u\in H_0^1(U)$, the Sobolev embedding gives
\begin{align*}
H_0^1(U)\hookrightarrow L^q(U),
\end{align*}
because $q<2<\frac{2n}{n-2}$ (since $n>3$).

Therefore,
\begin{align*}
\|u\|_{L^q(U)} \le C\|u\|_{H_0^1(U)}.
\end{align*}
Since on $H_0^1(U)$,
\begin{align*}
\|u\|_{H_0^1(U)} \sim \|\nabla u\|_{L^2(U)}
\end{align*}
(by Poincaré inequality), we get
\begin{align*}
\|u\|_{L^q(U)} \le C\|\nabla u\|_{L^2(U)}.
\end{align*}
Substitute into the previous inequality:
\begin{align*}
\|\nabla u\|_{L^2}^2 \le C\|\nabla u\|_{L^2}^q.
\end{align*}

---

\item \textbf{Step 4: Use that $q<2$}

Let
\begin{align*}
X=\|\nabla u\|_{L^2(U)}.
\end{align*}
Then
\begin{align*}
X^2\le CX^q.
\end{align*}
If $X>0$, divide by $X^q$:
\begin{align*}
X^{2-q}\le C.
\end{align*}
Thus,
\begin{align*}
X\le C^{1/(2-q)}.
\end{align*}
Hence,
\begin{align*}
\|\nabla u\|_{L^2(U)}\le C.
\end{align*}
By the equivalence of norms on $H_0^1(U)$,
\begin{align*}
\|u\|_{H_0^1(U)}\le C.
\end{align*}
This bound is independent of the weak solution (u).

---

\item \textbf{Final conclusion}

There exists $C>0$ such that every weak solution satisfies
\begin{align*}
\boxed{\|u\|_{H_0^1(U)}\le C.}
\end{align*}

---

\item \textbf{Main mechanism summarized}

We proved:
\begin{align*}
\|u\|_{H_0^1}^2 \lesssim \|u\|_{H_0^1}^q,
\end{align*}
and because $q<2$, this forces boundedness.

This is a standard \DEFINE{a priori estimate} argument in nonlinear Partial Differential Equations.

---
\end{description}

If you'd like, I can show how this estimate is used with the \DEFINE{Schaefer fixed point theorem} to prove existence of a weak solution for this nonlinear PDE.

\textbf{Questions:}
\begin{description}
	\item \textit{Why can we use `weak formulations`?}

	\item \textit{Explain Sobolev Embedding}.
\end{description}


%5%%%%%%%%%%%%%%%%%%%%
\newpage
\item Let $U$ be a bounded open subset of $\R^n$, $n>3$, let $p \in (2,2^*]$, and let 
\begin{align*}
	M = \BRACKET{u \in H_0^1(U)\,:\,\int_U |u|^p\,dx =1 }
\end{align*}
Consider the minimization problem
\begin{align*}
	c = \inf_{u\in M} \int_U |Du|^2\,dx
\end{align*}
\begin{enumerate}
	\item Show that $c > 0$.
	\item if $p \in (2, 2^*)$, show that the infimum is attained, i.e., $\exists u_0\in M$ such that
\begin{align*}
	c = \int_U |Du_0|^2\,dx
\end{align*}
\end{enumerate}

\HLINE

Let
\begin{align*}
M=\left\{u\in H_0^1(U):\int_U |u|^p\,dx=1\right\},
\end{align*}
and define
\begin{align*}
c=\inf_{u\in M}\int_U |\nabla u|^2\,dx.
\end{align*}
We want to prove:
\begin{enumerate}
\item $c>0$,
\item if $2<p<2^*$, then the infimum is attained.
\end{enumerate}
This is a standard application of the \DEFINE{Sobolev embedding theorem} and weak compactness in Calculus of Variations.

---
\begin{description}

\item \textbf{Part $1$: Show that $c>0$}

Take any $u\in M$. Since $u\in H_0^1(U)$, Sobolev embedding gives:
\begin{align*}
|u|*{L^p(U)}\le C|\nabla u|*{L^2(U)},
\end{align*}
because $p\le 2^*$.

Since $u\in M$,
\begin{align*}
|u|_{L^p(U)}=1.
\end{align*}
Therefore,
\begin{align*}
1\le C|\nabla u|_{L^2(U)}.
\end{align*}
Hence,
\begin{align*}
|\nabla u|_{L^2(U)}\ge \frac1C.
\end{align*}
Squaring,
\begin{align*}
\int_U |\nabla u|^2\,dx\ge \frac1{C^2}.
\end{align*}
Since this is true for every $u\in M$,
\begin{align*}
c\ge \frac1{C^2}>0.
\end{align*}
Thus,
\begin{align*}
\boxed{c>0.}
\end{align*}

---

\item \textbf{Part $2$: Show the infimum is attained for $2<p<2^*$}

We use the \DEFINE{direct method}.

---

	\begin{description}
    \item Step 1: Choose a minimizing sequence

Let $(u_j)\subset M$ be such that
\begin{align*}
\int_U |\nabla u_j|^2\,dx\to c.
\end{align*}
Then $(u_j)$ is bounded in $H_0^1(U)$, since
\begin{align*}
|\nabla u_j|_{L^2}^2\to c.
\end{align*}
Since on $H_0^1(U)$,
\begin{align*}
|u_j|*{H_0^1}\sim |\nabla u_j|*{L^2},
\end{align*}
the sequence is bounded in $H_0^1(U)$.

---

    \item Step 2: Extract a weakly convergent subsequence

Since $H_0^1(U)$ is reflexive, there exists a subsequence $still denoted (u_j$) and some $u_0\in H_0^1(U)$ such that

\begin{align*}
u_j\rightharpoonup u_0
\quad\text{weakly in }H_0^1(U).
\end{align*}
	\end{description}

---

    \item Step 3: Use compact Sobolev embedding

Since $2<p<2^*$, the embedding
\begin{align*}
H_0^1(U)\hookrightarrow L^p(U)
\end{align*}
is \DEFINE{compact}.

Therefore,
\begin{align*}
u_j\to u_0
\quad\text{strongly in }L^p(U).
\end{align*}
Hence,
\begin{align*}
\int_U |u_j|^p\,dx\to \int_U |u_0|^p\,dx.
\end{align*}
But each $u_j\in M$, so
\begin{align*}
\int_U |u_j|^p\,dx=1.
\end{align*}
Thus,
\begin{align*}
\int_U |u_0|^p\,dx=1.
\end{align*}
So
\begin{align*}
u_0\in M.
\end{align*}

---

    \item Step 4: Use weak lower semicontinuity

The Dirichlet energy is weakly lower semicontinuous:
\begin{align*}
\int_U |\nabla u_0|^2\,dx \le \liminf_{j\to\infty} \int_U |\nabla u_j|^2\,dx.
\end{align*}
Since the right side is $c$,
\begin{align*}
\int_U |\nabla u_0|^2\,dx\le c.
\end{align*}
But $u_0\in M$, so by definition of $c$,
\begin{align*}
c\le \int_U |\nabla u_0|^2\,dx.
\end{align*}
Therefore,
\begin{align*}
\int_U |\nabla u_0|^2\,dx=c.
\end{align*}
Thus the infimum is attained.

So there exists $u_0\in M$ such that
\begin{align*}
\boxed{c=\int_U |\nabla u_0|^2\,dx.}
\end{align*}

---

\item \textbf{Final conclusion}

		\begin{enumerate}
        \item
\begin{align*}
\boxed{c>0}
\end{align*}
by Sobolev inequality.

        \item For $2<p<2^*$, there exists $u_0\in M$ such that
\begin{align*}
\boxed{c=\int_U |\nabla u_0|^2\,dx}
\end{align*}
by:
\begin{enumerate}
\item bounded minimizing sequence,
\item weak compactness in $H_0^1(U)$,
\item compact embedding into $L^p(U)$,
\item weak lower semicontinuity.
\end{enumerate}
		\end{enumerate}

---
\end{description}

If you'd like, I can explain \DEFINE{why the argument fails at the critical exponent $p=2^*$}, which is where compactness breaks down.

\textbf{Questions:}
\begin{description}
\item \textit{Explain Calculus of Variations}
\item \textit{Explain semi-continuous}
\end{description}


%6%%%%%%%%%%%%%%%%%%%%
\newpage
\item Let $u \in C^2(\R\times(0,\infty))\cap C(\R\times[0,\infty))$ be the bouded solution of
\begin{align*}
	\left \{\begin{array}{rl}
		 u_t=u_{xx} & \text{in } \R\times(0,\infty)\\
		u = g & \text{on } \R\times\{0\}
	\end{array} \right .
\end{align*}
where $g \in C(\R)$ satisfies
\begin{align*}
	\lim_{x\to -\infty} g(x) = a, \quad \lim_{x\to +\infty} g(x) = b
\end{align*}
for some $a,b\in \R$.  Show that
\begin{align*}
	\lim_{x\to -\infty} u(x,t)=\frac{a+b}{2}\quad\forall x \in \R.
\end{align*}

\HLINE

Let
\begin{align*}
u_t=u_{xx}\qquad\text{in }\mathbb R\times (0,\infty),
\end{align*}
with initial data $u(x,0)=g(x)$, where $g$ is bounded and
\begin{align*}
\lim_{x\to -\infty}g(x)=a, \qquad \lim_{x\to +\infty}g(x)=b.
\end{align*}
We want to prove that for every $x\in\mathbb R$,
\begin{align*}
\lim_{t\to\infty}u(x,t)=\frac{a+b}{2}.
\end{align*}
(Your statement says $\lim_{x\to-\infty}u(x,t)$, but the conclusion $(a+b)/2$ is the \DEFINE{large time limit}, so the intended statement is almost certainly)
\begin{align*}
\boxed{\lim_{t\to\infty}u(x,t)=\frac{a+b}{2}}.
\end{align*}
\begin{description}

\item \textbf{Step 1: Write the solution using the heat kernel}

The bounded solution of the heat equation on $\mathbb R$ is given by
\begin{align*}
u(x,t)=\frac1{\sqrt{4\pi t}}
\int_{-\infty}^{\infty}
e^{-(x-y)^2/4t}g(y)\,dy.
\end{align*}
Equivalently, letting
\begin{align*}
z=\frac{x-y}{\sqrt{4t}},
\end{align*}
we get
\begin{align*}
u(x,t)=\frac1{\sqrt{\pi}}
\int_{-\infty}^{\infty}
e^{-z^2}g(x-2\sqrt{t},z)\,dz.
\end{align*}

\item \textbf{Step 2: Examine the behavior as $t\to\infty$}

Fix $x\in\mathbb R$.

Then for each fixed $z$,
\begin{align*}
x-2\sqrt{t},z\to
\begin{cases}
-\infty,& z>0,[1mm] \\
+\infty,& z<0.
\end{cases}
\end{align*}
Hence,
\begin{align*}
g(x-2\sqrt{t}z)\to
\begin{cases}
a,& z>0,[1mm] \\
b,& z<0.
\end{cases}
\end{align*}
Thus the integrand converges pointwise to
\begin{align*}
e^{-z^2}
\begin{cases}
a,& z>0,\\
b,& z<0.
\end{cases}
\end{align*}

\item \textbf{Step 3: Apply dominated convergence}

Since $g$ is bounded, say $|g(y)|\le M$, we have
\begin{align*}
\left|e^{-z^2}g(x-2\sqrt{t}z)\right| \le Me^{-z^2},
\end{align*}
and $Me^{-z^2}\in L^1(\mathbb R)$.

Therefore, by the dominated convergence theorem,
\begin{align*}
\lim_{t\to\infty}u(x,t) = \frac1{\sqrt{\pi}} \left( \int_0^\infty ae^{-z^2}\,dz + \int_{-\infty}^0 be^{-z^2}\,dz
\right).
\end{align*}

Factor out constants:
\begin{align*}
\frac1{\sqrt{\pi}} \left( a\int_0^\infty e^{-z^2}\,dz + b\int_{-\infty}^0 e^{-z^2}\,dz \right).
\end{align*}
But
\begin{align*}
\int_0^\infty e^{-z^2}\,dz = \left[\int_{-\infty}^0 e^{-z^2}\,dz] \frac{\sqrt{\pi}}2\right]
\end{align*}
Thus,
\begin{align*}
\lim_{t\to\infty}u(x,t) = \frac1{\sqrt{\pi}} \left( a\frac{\sqrt{\pi}}2+b\frac{\sqrt{\pi}}2 \right) = \frac{a+b}{2}.
\end{align*}
Therefore,
\begin{align*}
\boxed{
\lim_{t\to\infty}u(x,t)=\frac{a+b}{2} } \qquad\forall x\in\mathbb R.
\end{align*}

\item \textbf{Final conclusion}

The heat flow averages the left and right endstates of the initial data, and the solution converges pointwise to
\begin{align*}
\boxed{\frac{a+b}{2}}.
\end{align*}

This happens because the Gaussian kernel spreads symmetrically to the left and right as $t\to\infty$, assigning half the mass to each side.

This is a classic large-time asymptotic result for the heat equation.

\end{description}

If you'd like, I can explain the \DEFINE{intuition} for why the heat equation ``forgets" the shape of $g$ and only remembers the two end states $a$ and $b$.



%7%%%%%%%%%%%%%%%%%%%%
\newpage
\item Let $T>0$ and let $u \in C_1^2(\R^\times(0,T])\cap C(\R^n\times[0,T])$ solve
\begin{align*}
%	\left \{\begin{array}{rl}
%		 u_t-\Delta u=cu & \text{in } \R\times(0,T)\\
%		u = g & \text{on } \R\times\{0\}
%	\end{array} \right .
	\begin{cases}
		 u_t-\Delta u=cu & \text{in } \R\times(0,T)\\
		u = g & \text{on } \R\times\{0\}
	\end{cases}
\end{align*}
wher $c >0$ is a constant and $g\in C(\R^n)$ is bounded.  Assume that $u$ satisfies
\begin{align*}
	u(x,t) \le Ae^{a|x|^2}\quad \forall(x,t)\in\R^n\times[0,T]
\end{align*}
for some constants $A,a >0$.  Show that
\begin{align*}
	\sup_{\R^n\times[0,T]} u\le e^{cT}\sup_{\R^n} g
\end{align*}

\HLINE

We want to prove the estimate
\begin{align*}
\sup_{\mathbb R^n\times[0,T]}u \le e^{cT}\sup_{\mathbb R^n}g
\end{align*}
for a solution of
\begin{align*}
u_t-\Delta u=cu \qquad\text{in }\mathbb R^n\times(0,T),
\end{align*}
with initial data
\begin{align*}
u(x,0)=g(x),
\end{align*}
assuming
\begin{align*}
u(x,t)\le Ae^{a|x|^2}.
\end{align*}
This is a \DEFINE{maximum principle with a zeroth-order term}, standard in Partial Differential Equations.

\begin{description}

\item \textbf{Step 1: Remove the zeroth-order term}

Define
\begin{align*}
v(x,t)=e^{-ct}u(x,t).
\end{align*}
Then compute:
\begin{align*}
v_t = [-ce^{-ct}u+e^{-ct}u_t e^{-ct}(u_t-cu) ],
\end{align*}
and
\begin{align*}
\Delta v=e^{-ct}\Delta u.
\end{align*}
Thus,
\begin{align*}
v_t-\Delta v = e^{-ct}(u_t-\Delta u-cu).
\end{align*}
Since
\begin{align*}
u_t-\Delta u=cu,
\end{align*}
we obtain
\begin{align*}
v_t-\Delta v=0.
\end{align*}
So $v$ satisfies the heat equation:
\begin{align*}
v_t-\Delta v=0.
\end{align*}
Its initial data is
\begin{align*}
v(x,0)=u(x,0)=g(x).
\end{align*}

\item \textbf{Step 2: Apply the maximum principle}

We would like to conclude:
\begin{align*}
\sup_{\mathbb R^n\times[0,T]}v \le \sup_{\mathbb R^n}g.
\end{align*}
Because the domain is unbounded, we use the growth assumption
\begin{align*}
u(x,t)\le Ae^{a|x|^2}.
\end{align*}
Then
\begin{align*}
v(x,t)=e^{-ct}u(x,t)\le Ae^{a|x|^2},
\end{align*}
so $v$ has at most exponential growth.

This growth condition allows the maximum principle on $\mathbb R^n\times[0,T]$: bounded-above solutions of the heat equation attain no interior maximum above the initial maximum.

Hence,
\begin{align*}
\sup_{\mathbb R^n\times[0,T]}v \le \sup_{\mathbb R^n}g.
\end{align*}

\item \textbf{Step 3: Return to $u$}

Since
\begin{align*}
u(x,t)=e^{ct}v(x,t),
\end{align*}
we get
\begin{align*}
u(x,t) \le e^{ct}\sup_{\mathbb R^n}g \le e^{cT}\sup_{\mathbb R^n}g.
\end{align*}
Therefore,
\begin{align*}
\sup_{\mathbb R^n\times[0,T]}u \le e^{cT}\sup_{\mathbb R^n}g.
\end{align*}
Thus,
\begin{align*}
\boxed{
\sup_{\mathbb R^n\times[0,T]}u \le e^{cT}\sup_{\mathbb R^n}g
}
\end{align*}
which is what we wanted.

\item \textbf{Final idea}

The transformation
\begin{align*}
v=e^{-ct}u
\end{align*}
removes the reaction term $cu$, turning the PDE into the heat equation:
\begin{align*}
v_t-\Delta v=0.
\end{align*}
Then the \DEFINE{maximum principle} gives
\begin{align*}
v\le \sup g,
\end{align*}
and multiplying back gives
\begin{align*}
u\le e^{ct}\sup g\le e^{cT}\sup g.
\end{align*}
\end{description}

If you'd like, I can show how to justify the maximum principle on the \DEFINE{unbounded domain} $\mathbb R^n\times(0,T)$ using the barrier
\begin{align*}
v_\varepsilon(x,t)=v(x,t)-\varepsilon(|x|^2+\lambda t),
\end{align*}
which is the technical step behind this argument.



%8%%%%%%%%%%%%%%%%%%%%
\newpage
\item Let $u\in C^2(\R^2\times[0,\infty))$ be the solution of
\begin{align*}
	\left \{\begin{array}{rl}
		 u_{tt}= \Delta u & \text{in } \R\times(0,\infty)\\
		u = 0, u_t= h & \text{on } \R\times\{0\}
	\end{array} \right .
\end{align*}
where $g$ and $h$ are smooth with compact support.  Show that for each $x \in \R^2$, where exists a constant $C=C(x)>0$ such that 
\begin{align*}
	|u(x,t)| \le \frac{C(x)}{t}\quad\forall t\in (0,\infty).
\end{align*}

\HLINE

We are dealing with the \DEFINE{wave equation in two dimensions}

\begin{align*}
u_{tt}-\Delta u=0 \qquad\text{in }\mathbb R^2\times(0,\infty),
\end{align*}
with initial data
\begin{align*}
u(x,0)=0,\qquad u_t(x,0)=h(x),
\end{align*}
where $h\in C_c^\infty(\mathbb R^2)$.

We want to show that for each fixed $x\in\mathbb R^2$,
\begin{align*}
|u(x,t)|\le \frac{C(x)}{t}.
\end{align*}

This follows from the explicit solution formula for the 2D wave equation, known as \DEFINE{Poisson's formula}, a standard result in Partial Differential Equations.

\begin{description}
\item \textbf{Step 1: Use the 2D wave formula}

Since the initial displacement is zero, the solution is
\begin{align*}
u(x,t)=\frac1{2\pi}\int_{|x-y|<t} \frac{h(y)}{\sqrt{t^2-|x-y|^2}}\,dy.
\end{align*}
Because $h$ has compact support, there exists $R>0$ such that
\begin{align*}
\operatorname{supp}(h)\subset B_R(0).
\end{align*}
Thus the integral reduces to
\begin{align*}
u(x,t)=\frac1{2\pi} \int_{B_R(0)} \frac{h(y)}{\sqrt{t^2-|x-y|^2}}\,dy,
\end{align*}
whenever $t>|x-y|$ for $y\in B_R(0)$, i.e. for sufficiently large $t$.

\item \textbf{Step 2: Estimate the denominator}

Fix $x\in\mathbb R^2$.

For $y\in B_R(0$),
\begin{align*}
|x-y|\le |x|+R.
\end{align*}
Hence if $t>|x|+R$,
\begin{align*}
t^2-|x-y|^2\ge t^2-(|x|+R)^2.
\end{align*}
Therefore,
\begin{align*}
\frac1{\sqrt{t^2-|x-y|^2}} \le \frac1{\sqrt{t^2-(|x|+R)^2}}.
\end{align*}
Thus
\begin{align*}
|u(x,t)| \le \frac1{2\pi} \frac1{\sqrt{t^2-(|x|+R)^2}} \int_{B_R(0)}|h(y)|\,dy.
\end{align*}
Let
\begin{align*}
C_1=\frac1{2\pi}\int_{\mathbb R^2}|h(y)|\,dy.
\end{align*}
Then
\begin{align*}
|u(x,t)| \le \frac{C_1}{\sqrt{t^2-(|x|+R)^2}}.
\end{align*}

\item \textbf{Step 3: Compare with (1/t)}

For $t\ge 2(|x|+R)$,
\begin{align*}
t^2-(|x|+R)^2\ge \frac34 t^2.
\end{align*}
Hence
\begin{align*}
\sqrt{t^2-(|x|+R)^2}\ge \frac{\sqrt3}{2}t.
\end{align*}
So
\begin{align*}
|u(x,t)| \le \frac{2C_1}{\sqrt3}\frac1t.
\end{align*}
Thus for sufficiently large (t),
\begin{align*}
|u(x,t)|\le \frac{C}{t}.
\end{align*}

\item \textbf{Step 4: Extend to all (t>0)}

Since $u(x,t)$ is continuous in $t$, the function
\begin{align*}
t\mapsto t|u(x,t)|
\end{align*}
is continuous on $(0,2(|x|+R)]$, hence bounded there.

Thus there exists $C_2(x)$ such that
\begin{align*}
|u(x,t)|\le \frac{C_2(x)}{t}
\quad\text{for }0<t\le 2(|x|+R).
\end{align*}
Combining this with the large-(t) estimate gives a constant (C(x)) such that
\begin{align*}
|u(x,t)|\le \frac{C(x)}{t} \qquad\forall t>0.
\end{align*}
Hence,
\begin{align*}
\boxed{|u(x,t)|\le \frac{C(x)}{t}\quad\forall t>0.}
\end{align*}

\item \textbf{Final idea}

The solution is
\begin{align*}
u(x,t)=\frac1{2\pi}\int \frac{h(y)}{\sqrt{t^2-|x-y|^2}}\,dy,
\end{align*}
and because $h$ has compact support, the denominator behaves like $t$ for fixed $x$ as $t\to\infty$, giving
\begin{align*}
u(x,t)=O(t^{-1}).
\end{align*}
So the wave amplitude decays like (1/t) at each fixed point in two dimensions.

\end{description}

If you'd like, I can explain \DEFINE{why the decay is $1/t$ in dimension 2 but $1/t^{(n-1)/2}$ in higher dimensions}, which is the general dispersive decay law for the wave equation.



%9%%%%%%%%%%%%%%%%%%%%
\newpage
\item Consider the problem
\begin{align*}
	\left \{\begin{array}{rl}
		 u_{u} -\Delta u -2au_t+a^2u =f(x) & \text{in } \R\times(0,\infty)\\
		u = 0, u_t= 0 & \text{on } \R\times\{0\}
	\end{array} \right .
\end{align*}
where $a > 0$ is a constant and $f$ is a continous function on $\R^n$.  By considering the energy
\begin{align*}
	E(t) = \frac{1}{2} \int_{\R^n} (u_t^2+|\nabla u|^2 + a^23u^2)\,dx,
\end{align*}
show that this problem has at most one solution in $C_c^2(\R^n\times[0,\infty))$.

\HLINE


To prove \DEFINE{uniqueness}, we show that if $u$ is a solution of the homogeneous problem

\begin{align*}
\begin{cases}
u_{tt}-\Delta u-2au_t+a^2u=0 & \text{in }\mathbb R^n\times (0,\infty),[1mm]
u(x,0)=0,\quad u_t(x,0)=0,
\end{cases}
\end{align*}
then $u\equiv 0$.

This implies the original problem has \DEFINE{at most one solution}, since the difference of two solutions satisfies this homogeneous problem.

This is a standard \DEFINE{energy method} argument in Partial Differential Equations.

\begin{description}

\item \textbf{Step 1: Define the energy}

Let
\begin{align*}
E(t)=\frac12\int_{\mathbb R^n}\left(u_t^2+|\nabla u|^2+a^2u^2\right)\,dx.
\end{align*}

We compute $E'(t$).

\item \textbf{Step 2: Differentiate the energy}

Differentiate under the integral:

\begin{align*}
E'(t) = \int_{\mathbb R^n}\left(u_tu_{tt}+\nabla u\cdot\nabla u_t+a^2uu_t\right)\,dx.
\end{align*}

Integrate the middle term by parts:

\begin{align*}
\int_{\mathbb R^n}\nabla u\cdot\nabla u_t\,dx = -\int_{\mathbb R^n}(\Delta u)u_t\,dx,
\end{align*}

since $u\in C_c^2$, so there are no boundary terms.

Thus,
\begin{align*}
E'(t) = \int_{\mathbb R^n}u_t(u_{tt}-\Delta u+a^2u)\,dx.
\end{align*}

\item \textbf{Step 3: Use the PDE}

The PDE is
\begin{align*}
u_{tt}-\Delta u-2au_t+a^2u=0,
\end{align*}
so
\begin{align*}
u_{tt}-\Delta u+a^2u=2au_t.
\end{align*}
Substitute:
\begin{align*}
E'(t) = \int_{\mathbb R^n}u_t(2au_t)\,dx 2a\int_{\mathbb R^n}u_t^2\,dx.
\end{align*}
Therefore,
\begin{align*}
E'(t)=2a\int_{\mathbb R^n}u_t^2\,dx\ge 0.
\end{align*}
So the energy is nondecreasing.

\item \textbf{Step 4: Use the initial conditions}

Because
\begin{align*}
u(x,0)=0,\qquad u_t(x,0)=0,
\end{align*}
we have
\begin{align*}
E(0)=0.
\end{align*}
Since $E'(t)\ge 0$,
\begin{align*}
E(t)\ge 0.
\end{align*}
But this alone does not give uniqueness—we need a better quantity.

\item \textbf{Step 5: Introduce the transformed variable}

Define
\begin{align*}
v(x,t)=e^{-at}u(x,t).
\end{align*}
Then compute:
\begin{align*}
u&=e^{at}v,\\
u_t&=e^{at}(v_t+av),\\
u_{tt}&=e^{at}(v_{tt}+2av_t+a^2v).
\end{align*}
Substitute into
\begin{align*}
u_{tt}-\Delta u-2au_t+a^2u&=0:\\
e^{at}(v_{tt}+2av_t+a^2v)-e^{at}\Delta v-2ae^{at}(v_t+av)+a^2e^{at}v&=0.
\end{align*}
Simplify:
\begin{align*}
e^{at}(v_{tt}-\Delta v)=0.
\end{align*}
Hence $v$ solves the wave equation
\begin{align*}
v_{tt}-\Delta v=0.
\end{align*}
Also,
\begin{align*}
v(x,0)=u(x,0)=0,
\end{align*}
and
\begin{align*}
v_t(x,0)=u_t(x,0)-au(x,0)=0.
\end{align*}
Thus $v$ satisfies
\begin{align*}
\begin{cases}
v_{tt}-\Delta v=0,\
v(x,0)=0,\quad v_t(x,0)=0.
\end{cases}
\end{align*}

\item \textbf{Step 6: Apply the standard wave energy identity}

Define
\begin{align*}
F(t)=\frac12\int_{\mathbb R^n}(v_t^2+|\nabla v|^2)\,dx.
\end{align*}
For the wave equation,
\begin{align*}
F'(t)=0.
\end{align*}
Since $F(0)=0$, we get
\begin{align*}
F(t)=0.
\end{align*}
Therefore,
\begin{align*}
v_t=0,\qquad \nabla v=0.
\end{align*}
So $v\equiv 0$.
Hence
\begin{align*}
u=e^{at}v\equiv 0.
\end{align*}
Thus the homogeneous problem has only the zero solution.

Therefore the original problem has at most one solution.

\item \textbf{Final conclusion}

The difference of any two solutions satisfies the homogeneous equation, and the transformation

\begin{align*}
v=e^{-at}u
\end{align*}
reduces it to the classical wave equation with zero initial data, whose energy remains zero.

Hence
\begin{align*}
u\equiv 0.
\end{align*}
Therefore the original problem has at most one solution:
\begin{align*}
\boxed{\text{The problem has at most one solution in }C_c^2(\mathbb R^n\times[0,\infty)).}
\end{align*}

\end{description}

If you'd like, I can show how to prove uniqueness \DEFINE{directly from the given energy $E(t)$} without introducing the substitution $v=e^{-at}u$, since that is probably what the exercise intends.



%10%%%%%%%%%%%%%%%%%%%
\newpage
\item Show that the problem
\begin{align*}
	\left \{\begin{array}{rl}
		 u_{u} -\Delta u  =f(x) & \text{in } U\times(0,T)\\
		u = g, u_t= h & \text{on } U\times\{0\} \\
		u_t + \PART{u}{\nu} = 0 & \text{on } \partial U \times[0,T],
	\end{array} \right .
\end{align*}
where $U\subset \R^n$ is open and bounded with smooth boundary $\partial U, T>0$, and $\nu$ is the outward unit normal to $\partial U$, has at most one solution $u \in C^2(\CONJ{U}\times[0,T])$.

\end{enumerate}

\HLINE


To show the problem has \DEFINE{at most one solution}, we prove that if $u_1,u_2$ are two solutions, then they must be equal.

The standard method is an \DEFINE{energy argument}.

This is a uniqueness proof for the wave equation with \DEFINE{dissipative boundary conditions}, a standard result in Partial Differential Equations.

\begin{description}

\item \textbf{Step 1: Take the difference of two solutions}

Suppose $u_1,u_2\in C^2(\overline U\times[0,T])$ are two solutions.

Let
\begin{align*}
w=u_1-u_2.
\end{align*}
Then $w$ satisfies the homogeneous problem:
\begin{align*}
\begin{cases}
w_{tt}-\Delta w=0 &\text{in }U\times(0,T),\\
w(x,0)=0,\quad w_t(x,0)=0 &\text{in }U,\\
w_t+\dfrac{\partial w}{\partial \nu}=0 &\text{on }\partial U\times[0,T].
\end{cases}
\end{align*}
We will show that $w\equiv 0$.

\item \textbf{Step 2: Define the energy}

Define
\begin{align*}
E(t)=\frac12\int_U \left(w_t^2+|\nabla w|^2\right)\,dx.
\end{align*}
This is the usual wave energy.

\item \textbf{Step 3: Differentiate the energy}

Differentiate:
\begin{align*}
E'(t)=\int_U w_tw_{tt}\,dx+\int_U \nabla w\cdot \nabla w_t\,dx.
\end{align*}
Integrate the second term by parts:
\begin{align*}
\int_U \nabla w\cdot\nabla w_t\,dx = -\int_U (\Delta w)w_t\,dx +\int_{\partial U}\frac{\partial w}{\partial \nu}w_t\,dS.
\end{align*}
Thus,
\begin{align*}
E'(t) = \int_U w_t(w_{tt}-\Delta w)\,dx +\int_{\partial U}\frac{\partial w}{\partial \nu}w_t\,dS.
\end{align*}
Since $w_{tt}-\Delta w=0$,
\begin{align*}
E'(t) = \int_{\partial U}\frac{\partial w}{\partial \nu}w_t\,dS.
\end{align*}

\item \textbf{Step 4: Use the boundary condition}

On $\partial U$,
\begin{align*}
w_t+\frac{\partial w}{\partial \nu}=0,
\end{align*}
so
\begin{align*}
\frac{\partial w}{\partial \nu}=-w_t.
\end{align*}
Substitute:
\begin{align*}
E'(t) = \int_{\partial U}(-w_t)w_t\,dS = -\int_{\partial U}w_t^2\,dS.
\end{align*}
Hence,
\begin{align*}
E'(t)\le 0.
\end{align*}
So the energy is nonincreasing.

\item \textbf{Step 5: Use the initial conditions}

Because
\begin{align*}
w(x,0)=0,\qquad w_t(x,0)=0,
\end{align*}
we have
\begin{align*}
E(0)=0.
\end{align*}
Since $E(t)\ge 0$ and $E'(t)\le 0$,
\begin{align*}
E(t)\equiv 0.
\end{align*}
Thus,
\begin{align*}
\int_U \left(w_t^2+|\nabla w|^2\right)\,dx=0.
\end{align*}
Hence,
\begin{align*}
w_t=0,\qquad \nabla w=0
\end{align*}
in $U\times(0,T)$.
So $w$ is constant in space and time.

Since $w(x,0)=0$,
\begin{align*}
w\equiv 0.
\end{align*}
Thus,
\begin{align*}
u_1=u_2.
\end{align*}

\item \textbf{Final conclusion}

The difference of two solutions has zero energy, hence is identically zero.

Therefore the problem has at most one solution:
\begin{align*}
\boxed{\text{There is at most one }u\in C^2(\overline U\times[0,T])\text{ solving the problem.}}
\end{align*}

\item \textbf{Key idea}

The dissipative boundary condition
\begin{align*}
u_t+\frac{\partial u}{\partial \nu}=0
\end{align*}
forces
\begin{align*}
E'(t)=-\int_{\partial U}u_t^2\,dS\le 0,
\end{align*}
so the energy decreases, and zero initial energy implies the solution is zero.

This is the fundamental uniqueness mechanism for damped boundary wave problems.

\end{description}

If you'd like, I can explain \DEFINE{why the boundary condition $u_t+\partial_\nu u=0$ is called absorbing/dissipative}, which gives intuition for why the energy decreases.

\textbf{Questions}

\begin{description}
	\item \textit{Explain the `energy argument'.}
	\item \textit{What are the energy formulas for the standard PDEs}.
\end{description}


%%%%%%%%%%%%%%%%%%%% Notes
\newpage
\section*{Notes}

\begin{description}
	\item \DEFINE{Explain `weak convergence'}

\DEFINE{``Weakly convergent"} means a sequence converges \DEFINE{when tested against every continuous linear functional}, even if it does \DEFINE{not} converge in norm.

This is weaker than ordinary (strong) convergence.

---
\begin{description}

\item \textbf{1. Definition of weak convergence}

Let $X$ be a Banach space.

A sequence $u_j$ \DEFINE{converges weakly} to $u$ in $X$, written
\begin{align*}
u_j \rightharpoonup u,
\end{align*}
if for every bounded linear functional $L\in X^*$,
\begin{align*}
L(u_j)\to L(u).
\end{align*}

---

\item \textbf{2. Intuition}

Instead of requiring
\begin{align*}
|u_j-u|_X\to 0
\end{align*}
(which is \DEFINE{strong convergence}), weak convergence only requires that
\begin{align*}
L(u_j-u)\to 0
\end{align*}

for every linear measurement $L$.

So:
\begin{itemize}
\item \DEFINE{strong convergence:} vectors themselves get close;
\item \DEFINE{weak convergence:} all linear “probes” of the vectors converge.
\end{itemize}

---

\item \textbf{3. Example in $L^2(U)$}

A sequence $f_j$ converges weakly to $f$ in $L^2(U)$ if
\begin{align*}
\int_U f_j g\,dx \to \int_U fg\,dx
\qquad\forall g\in L^2(U).
\end{align*}
So instead of
\begin{align*}
|f_j-f|_{L^2}\to 0,
\end{align*}

we only require convergence of all inner products.

Because $L^2(U)$ is a Hilbert space, weak convergence is:
\begin{align*}
f_j\rightharpoonup f \iff \langle f_j,g\rangle\to\langle f,g\rangle \quad\forall g\in L^2(U).
\end{align*}

---

\item \textbf{4. Example in $H_0^1(U)$}

In $H_0^1(U)$,
\begin{align*}
u_j\rightharpoonup u
\end{align*}
means:
\begin{align*}
\int_U \nabla u_j\cdot\nabla v\,dx \to \int_U \nabla u\cdot\nabla v\,dx \qquad \forall v\in H_0^1(U),
\end{align*}
and similarly for the $L^2$-part.

So weak convergence in Sobolev spaces means the derivatives converge “in the averaged sense.”

---

\item \textbf{5. Strong convergence implies weak convergence}

If
\begin{align*}
u_j\to u
\quad\text{strongly in }X,
\end{align*}
then automatically
\begin{align*}
u_j\rightharpoonup u.
\end{align*}
Because if $L$ is continuous,
\begin{align*}
|L(u_j)-L(u)| \le \|L\|\,\|u_j-u\|,
\end{align*}
so the right-hand side goes to zero.

Thus:
\begin{align*}
\text{strong convergence } \implies \text{ weak convergence}.
\end{align*}
But the converse is generally false.

---

\item \textbf{6. Weak but not strong convergence example}

In $L^2(0,2\pi)$, define
\begin{align*}
f_j(x)=\sin(jx).
\end{align*}
Then:
\begin{description}
\item $f_j$ does \DEFINE{not} converge strongly to $0$, because
\begin{align*}
  |f_j|_{L^2}\not\to 0,
\end{align*}
but
\item $f_j\rightharpoonup 0$ weakly, because
\begin{align*}
  \int_0^{2\pi}\sin(jx)g(x)\,dx\to 0
\end{align*}
for every $g\in L^2$.
\end{description}

So the oscillations average out.

This is the intuition:

\> weak convergence allows oscillation, as long as averages converge.

---

\item \textbf{7. Why weak convergence matters}

Weak convergence is crucial because bounded sequences may fail to converge strongly, but often have weakly convergent subsequences.

This is the idea behind:

\> \DEFINE{bounded sequence + reflexive space $\to$ weakly convergent subsequence}

For example:
\begin{itemize}
\item bounded sequences in $L^p$, $1<p<\infty$,
\item bounded sequences in $H^1$,
\end{itemize}
have weakly convergent subsequences.

This is one of the main compactness tools in Functional Analysis.

---

\item \textbf{8. In your PDE problem}

You had:
\begin{align*}
u_j \rightharpoonup u \quad\text{in } H_0^1(U),
\end{align*}
which means:
\begin{itemize}
\item $u_j$ may not converge in norm,
\item but for every test function $v$,
\end{itemize}

\begin{align*}
\int_U \nabla u_j\cdot\nabla v\,dx
\to
\int_U \nabla u\cdot\nabla v\,dx.
\end{align*}
That is exactly what lets us pass to the limit in the weak formulation.

---

\item \textbf{Final summary}

Weak convergence means:
\begin{align*}
u_j\rightharpoonup u
\end{align*}
if all linear functionals converge:
\begin{align*}
L(u_j)\to L(u).
\end{align*}

It is weaker than norm convergence, but bounded sequences often have weakly convergent subsequences, which makes it fundamental in PDEs and Sobolev spaces.

---
\end{description}

If you'd like, I can explain \DEFINE{why bounded sequences in $H_0^1(U)$ always have weakly convergent subsequences}, which is the compactness fact used in your proof.

\HLINE

%%%%%%%%%%%%%%%%%%%%%%%
	\newpage
	\item \DEFINE{Explain `weak differentiation'}

A \DEFINE{weak derivative} is a generalized notion of derivative that lets us differentiate functions that may not be classically differentiable.

It is fundamental in Partial Differential Equations because weak solutions of PDEs are defined using weak derivatives.

---
\begin{description}

\item \textbf{1. Motivation}

In classical calculus, the derivative $u'(x)$ is defined pointwise:
\begin{align*}
u'(x)=\lim_{h\to 0}\frac{u(x+h)-u(x)}{h}.
\end{align*}
But many functions:
\begin{itemize}
\item are not differentiable everywhere,
\item may have corners,
\item may only be integrable,
\end{itemize}
yet we still want to talk about their derivatives.

The weak derivative does this by moving the derivative onto a \DEFINE{test function} using integration by parts.

---

\item \textbf{2. Classical integration by parts idea}

If $u,v\in C^1(a,b)$, then
\begin{align*}
\int_a^b u'(x)v(x)\,dx = -\int_a^b u(x)v'(x)\,dx
\end{align*}
assuming $v(a)=v(b)=0$.

Notice:

\> This identity characterizes $u'$ without using pointwise differentiation.

That motivates the definition.

---

\item \textbf{3. Definition of weak derivative in one dimension}

Let $u\in L^1_{\mathrm{loc}}(a,b)$.

A function $g\in L^1_{\mathrm{loc}}(a,b)$ is the \DEFINE{weak derivative} of $u$ if
\begin{align*}
\int_a^b u(x)\varphi'(x)\,dx = -\int_a^b g(x)\varphi(x)\,dx 
\end{align*}
\begin{align*}
\varphi\in C_c^\infty(a,b).
\end{align*}
Then we write
\begin{align*}
g=u'.
\end{align*}

---

\item \textbf{4. Interpretation}

Instead of requiring:
\begin{align*}
u'(x)=g(x)
\end{align*}
pointwise, we require:
\begin{align*}
\int u\,\varphi'=-\int g\,\varphi
\end{align*}
for every smooth compactly supported $\varphi$.

So:

\> \DEFINE{$g$ acts like the derivative of $u$ under integration.}

---

\item \textbf{5. Example: $u(x)=|x|$}

Classically:
\begin{align*}
u(x)=|x|
\end{align*}
is not differentiable at $x=0$.

But define
\begin{align*}
g(x)= \begin{cases}
1,&x>0,\\
-1,&x<0.
\end{cases}
\end{align*}
Then for every $\varphi\in C_c^\infty(\mathbb R)$,
\begin{align*}
\int_{\mathbb R}|x|\varphi'(x),dx = -\int_{\mathbb R}g(x)\varphi(x),dx.
\end{align*}
So $g$ is the weak derivative of $|x|$.

Thus:
\begin{align*}
(|x|)'=\operatorname{sgn}(x)
\end{align*}

in the weak sense.

This is valid even though the classical derivative fails at $0$.

---

\item \textbf{6. Multivariable weak derivatives}

For $u\in L^1_{\text{loc}}(U)$, the weak partial derivative $\partial u/\partial x_i$ is a function $g\in L^1_{\text{loc}}(U)$ such that

\begin{align*}
\int_U u,\frac{\partial \varphi}{\partial x_i},dx = -\int_U g,\varphi,dx
\end{align*}
for all $\varphi\in C_c^\infty(U)$.

Then we write
\begin{align*}
g=\frac{\partial u}{\partial x_i}.
\end{align*}

---

\item \textbf{7. Why weak derivatives matter}

Many PDE solutions are not smooth enough for classical derivatives.

Example:
\begin{align*}
-\Delta u=f
\end{align*}
may not have a $C^2$ solution, but may have a solution with
\begin{align*}
u\in H^1(U),
\end{align*}
meaning:
\begin{itemize}
\item $u\in L^2(U)$,
\item weak first derivatives are in $L^2(U)$.
\end{itemize}

Thus weak derivatives allow us to define Sobolev spaces:
\begin{align*}
W^{k,p}(U) = {u\in L^p(U): D^\alpha u\in L^p(U)}.
\end{align*}

---

\item \textbf{8. If classical derivative exists, weak derivative agrees}

If $u\in C^1(U)$, then its weak derivative equals its classical derivative.

So weak derivatives extend the classical notion—they do not replace it.

---

\item \textbf{9. Key intuition}

A weak derivative is a derivative defined \DEFINE{through integration against test functions}, rather than pointwise limits.

This lets us differentiate rough functions.

---

\item \textbf{Final summary}

A function $g$ is the weak derivative of $u$ if
\begin{align*}
\int_U u,D\varphi,dx = -\int_U g,\varphi,dx
\end{align*}
for every smooth compactly supported test function $\varphi$.

This is the generalized derivative used in Sobolev spaces and weak PDE solutions.

---
\end{description}

If you'd like, I can show next \DEFINE{why the weak derivative is unique}, which is an important theorem behind the definition of Sobolev spaces.


%%%%%%%%%%%%%%%%%%%%%%%
	\newpage
	\item \DEFINE{Explain the relationship between $H_0^1$ and $L^p$ spaces}

The relationship between $H_0^1(U)$ and $L^p(U)$ spaces comes from the fact that functions in $H_0^1(U)$ automatically belong to certain $L^p(U)$ spaces. This is given by the \DEFINE{Sobolev embedding theorem}, one of the central results in Functional Analysis.

---
\begin{description}

\item \textbf{1. Recall the definitions}

---

\begin{itemize}
\item \textbf{$L^p(U)$}

A function $u$ is in $L^p(U)$ if
\begin{align*}
\int_U |u|^p,dx<\infty.
\end{align*}

This measures the \DEFINE{size of the function}.

---

\item \textbf{$H^1(U)$}

\begin{align*}
H^1(U)=W^{1,2}(U)
\end{align*}
means:
\begin{align*}
u\in L^2(U)
\quad\text{and}\quad
\nabla u\in L^2(U).
\end{align*}

So $H^1$ measures:

\begin{enumerate}
\item the size of $u$,
\item the size of its gradient.
\end{enumerate}

---

\item \textbf{$H_0^1(U)$}

$H_0^1(U)$ is the closure of $C_c^\infty(U)$ in the $H^1$-norm.

These are $H^1$ functions whose boundary trace is zero.

---
\end{itemize}

\item \textbf{2. $H_0^1(U)\subset L^p(U)$}

If $U\subset \mathbb R^n$ is bounded, then every function in $H_0^1(U)$ lies in $L^p(U)$ for a range of $p$.

Specifically:

---
\begin{itemize}

\item \textbf{If $n>2$,}
\begin{align*}
H_0^1(U)\hookrightarrow L^p(U)
\quad\text{for }1\le p\le \frac{2n}{n-2}.
\end{align*}
That means:
\begin{align*}
u\in H_0^1(U)\implies u\in L^p(U),
\end{align*}
and there exists $C>0$ such that
\begin{align*}
\|u\|_{L^p(U)} \le C\|u\|_{H_0^1(U)}.
\end{align*}
The exponent
\begin{align*}
2^*=\frac{2n}{n-2}
\end{align*}
is called the \DEFINE{critical Sobolev exponent}.

---

\item \textbf{If $n=2$,}
\begin{align*}
H_0^1(U)\hookrightarrow L^p(U)
\quad\forall p<\infty.
\end{align*}
So every $H_0^1$ function is in every finite $L^p$ space.

---

\item \textbf{If $n=1$,}
\begin{align*}
H_0^1(U)\hookrightarrow L^\infty(U).
\end{align*}
So one-dimensional $H_0^1$ functions are bounded.

---
\end{itemize}

\item \textbf{3. Why this matters}

The space $H_0^1(U)$ gives you \DEFINE{extra integrability}.

Knowing only $u\in L^2$ gives:
\begin{align*}
u\in L^2.
\end{align*}
But knowing $u\in H_0^1$ gives:
\begin{align*}
u\in L^p
\end{align*}
for larger (p).

So:

\> the derivative information improves the integrability of the function.

That is the key relationship.

---

\item \textbf{4. Example in dimension $3$}

If $n=3$,
\begin{align*}
2^*=\frac{2(3)}{3-2}=6.
\end{align*}
So
\begin{align*}
H_0^1(U)\hookrightarrow L^6(U).
\end{align*}
Thus every $u\in H_0^1(U)$ satisfies
\begin{align*}
\|u\|_{L^6}\le C\|\nabla u\|_{L^2}.
\end{align*}
This is often written as
\begin{align*}
H_0^1(U)\subset L^6(U).
\end{align*}
This is essential in nonlinear elliptic PDEs.

---

\item \textbf{5. Poincaré inequality connects them}

For $u\in H_0^1(U)$,
\begin{align*}
\|u\|_{L^2(U)} \le C\|\nabla u\|_{L^2(U)}.
\end{align*}
This means that in $H_0^1(U)$, the gradient controls the $L^2$-norm.

Therefore,
\begin{align*}
\|u\|_{H_0^1(U)} \sim \|\nabla u\|_{L^2(U)}.
\end{align*}
Combining this with Sobolev embedding gives:
\begin{align*}
\|u\|_{L^p(U)} \le C\|\nabla u\|_{L^2(U)}.
\end{align*}
This is the fundamental bridge between $H_0^1$ and $L^p$.

---

\item \textbf{6. Compact embeddings}

The embedding
\begin{align*}
H_0^1(U)\hookrightarrow L^p(U)
\end{align*}
is not only continuous, but compact for
\begin{align*}
1\le p<\frac{2n}{n-2}.
\end{align*}
That means:

If $(u_j)$ is bounded in $H_0^1(U)$, then some subsequence converges \DEFINE{strongly} in $L^p(U)$.

This is the \DEFINE{Rellich–Kondrachov compactness theorem}, fundamental in Partial Differential Equations.

---

\item \textbf{7. Summary of the relationship}

The space $H_0^1(U)$ is stronger than $L^p(U)$:
\begin{align*}
H_0^1(U)\subset L^p(U)
\end{align*}
for appropriate $p$.

The inclusion is controlled by Sobolev inequalities:
\begin{align*}
\|u\|_{L^p}\le C\|u\|_{H_0^1}.
\end{align*}
So:

\> $H_0^1(U$) gives both square integrability \DEFINE{and} derivative control, and that forces better $L^p$ behavior.

---

\item \textbf{Final summary}

For bounded (U),
\begin{align*}
H_0^1(U)\hookrightarrow L^p(U)
\end{align*}
with:
\begin{align*}
1\le p\le \frac{2n}{n-2}\quad (n>2).
\end{align*}

This means every $H_0^1$ function automatically lies in certain $L^p$ spaces, and its $L^p$-norm is controlled by its $H_0^1$-norm.

That relationship is one of the core structural facts in Sobolev space theory.

---
\end{description}

If you'd like, next I can explain \DEFINE{why the exponent $\frac{2n}{n-2}$ appears}, since that critical exponent is extremely important in PDE.


\newpage

	\item \DEFINE{Why can we use `weak formulations`?}

We use \DEFINE{weak formulations} because many differential equations do not have sufficiently smooth \DEFINE{classical solutions}, but they often do have solutions in a weaker sense.

The weak formulation allows us to:
\begin{enumerate}
\item \DEFINE{make sense of derivatives for rough functions},
\item \DEFINE{prove existence of solutions},
\item \DEFINE{work in Sobolev spaces},
\item \DEFINE{use functional analysis tools like Lax–Milgram}.
\end{enumerate}
This is foundational in Partial Differential Equations.

---
\begin{description}

\item \textbf{1. Classical solutions may not exist}

Take the PDE
\begin{align*}
-\Delta u=f.
\end{align*}
A \DEFINE{classical solution} requires
\begin{align*}
u\in C^2(U),
\end{align*}
because $\Delta u$ involves second derivatives.

But even if $f$ is nice, the solution may fail to be $C^2$.

For example:
\begin{itemize}
\item the domain may be irregular,
\item the coefficients may be rough,
\item the forcing term $f$ may only be in $L^2$.
\end{itemize}
Then classical derivatives may not exist.

---

\item \textbf{2. Weak formulations lower the differentiability requirement}

Suppose
\begin{align*}
-\Delta u=f.
\end{align*}
Multiply by a test function $v\in C_c^\infty(U)$:
\begin{align*}
\int_U (-\Delta u)v,dx=\int_U fv,dx.
\end{align*}
Integrate by parts:
\begin{align*}
\int_U \nabla u\cdot\nabla v,dx=\int_U fv,dx.
\end{align*}
Now only \DEFINE{first derivatives} of $u$ appear.

So instead of requiring:
\begin{align*}
u\in C^2(U),
\end{align*}
we only need:
\begin{align*}
u\in H^1(U).
\end{align*}
That is a much larger space.

So:

\> the weak formulation reduces smoothness requirements.

---

\item \textbf{3. Weak derivatives let us define derivatives for rough functions}

In Sobolev spaces, derivatives are interpreted weakly:
\begin{align*}
\int_U u,D\varphi = -\int_U (Du)\varphi.
\end{align*}
This means $u$ need not be classically differentiable.

Example:
\begin{align*}
u(x)=|x|
\end{align*}
is not differentiable at $0$, but it has a weak derivative.

Thus weak formulations let us solve PDEs with non-smooth solutions.

---

\item \textbf{4. Weak formulations turn PDEs into integral equations}

Instead of solving
\begin{align*}
-\Delta u=f,
\end{align*}
we solve:
\begin{align*}
\int_U \nabla u\cdot\nabla v,dx = \int_U fv,dx \qquad \forall v.
\end{align*}
This is an equation in a Hilbert space.

That allows the use of tools from Functional Analysis, such as:
\begin{itemize}
\item Lax–Milgram,
\item weak compactness,
\item variational methods,
\item monotone operator methods.
\end{itemize}
These tools are often unavailable in the classical setting.

---

\item \textbf{5. Weak formulations are stable under limits}

Suppose $u_j$ solves approximating equations and
\begin{align*}
u_j\rightharpoonup u.
\end{align*}
The weak formulation:
\begin{align*}
\int_U \nabla u_j\cdot\nabla v,dx=\int_U f_jv,dx
\end{align*}
passes to the limit easily.

That is because integrals behave well under weak convergence.

But pointwise second derivatives generally do \DEFINE{not} pass to limits.

So:

\> weak formulations are compatible with approximation.

This is essential in existence proofs.

---

\item \textbf{6. Physical meaning}

Weak formulations often represent the \DEFINE{integral balance law} of the PDE.

For example, in elasticity or heat flow, the weak form expresses:
\begin{itemize}
\item conservation of energy,
\item conservation of mass,
\item equilibrium.
\end{itemize}
Thus weak solutions are often physically meaningful even when classical derivatives fail.

---

\item \textbf{7. Classical solutions imply weak solutions}
If (u) is smooth, then integration by parts shows it satisfies the weak formulation.

So:

\> every classical solution is a weak solution.

That means weak solutions generalize classical ones.

Moreover, if the weak solution is smooth enough, it satisfies the PDE classically.

So the weak formulation is not changing the equation—it is extending the notion of solution.

---

\item \textbf{Final summary}

We use weak formulations because they:
\begin{enumerate}
\item allow less regular solutions,
\item reduce derivative requirements,
\item convert PDEs into integral equations,
\item enable functional analytic methods,
\item are stable under weak convergence.
\end{enumerate}

In short:

\> \DEFINE{Weak formulations let us solve PDEs in spaces where classical derivatives may not exist.}

That is why they are central in modern PDE theory.

---
\end{description}

If you'd like, I can show the exact derivation of the weak formulation for
\begin{align*}
-\Delta u=f,\quad u=0,
\end{align*}
step by step using integration by parts.

%%%%%%%%%%%%%%%%%%5
\newpage
	\item \DEFINE{Explain Sobolev Embedding}.

The \DEFINE{Sobolev Embedding Theorem} tells us that if a function has enough integrable derivatives, then the function itself automatically belongs to a “better” function space—typically an $L^p$ space or even a continuous function space.

It is one of the central structural theorems in Functional Analysis and underlies much of modern Partial Differential Equations.

---
\begin{description}

\item \textbf{1. The basic idea}

Sobolev spaces measure both:

\begin{enumerate}
\item the size of the function,
\item the size of its derivatives.
\end{enumerate}
The Sobolev embedding theorem says:

\> \DEFINE{Derivative control improves integrability (and sometimes continuity).}

For example:
\begin{align*}
u\in H_0^1(U)
\end{align*}
implies
\begin{align*}
u\in L^p(U)
\end{align*}
for certain $p>2$.

So knowing $\nabla u\in L^2$ gives stronger information about $u$.

\item \textbf{2. The classical Sobolev embedding}

Let $U\subset \mathbb R^n$ be bounded.

If $n>2$, then
\begin{align*}
H_0^1(U)\hookrightarrow L^{2^*}(U),
\end{align*}
where
\begin{align*}
2^*=\frac{2n}{n-2}.
\end{align*}
This means:
\begin{enumerate}
\item every $u\in H_0^1(U)$ belongs to $L^{2^*}(U)$,
\item there exists $C>0$ such that
\end{enumerate}
\begin{align*}
\|u\|_{L^{2^*}(U)} \le C\|u\|_{H_0^1(U)}.
\end{align*}
Because of Poincaré inequality, this is often written as
\begin{align*}
\|u\|_{L^{2^*}(U)} \le C\|\nabla u\|_{L^2(U)}.
\end{align*}
This is the fundamental Sobolev inequality.

\item \textbf{3. Meaning of the embedding notation}

The notation
\begin{align*}
H_0^1(U)\hookrightarrow L^p(U)
\end{align*}
means:
\begin{itemize}
\item every $u\in H_0^1(U)$ is in $L^p(U)$,
\item the inclusion map is continuous.
\end{itemize}
That is:
\begin{align*}
\|u\|_{L^p(U)} \le C\|u\|_{H_0^1(U)}.
\end{align*}
So the Sobolev norm controls the $L^p$-norm.

\item \textbf{4. The critical exponent}

The number
\begin{align*}
2^*=\frac{2n}{n-2}
\end{align*}
is called the \DEFINE{critical Sobolev exponent}.

The theorem says:
\begin{align*}
H_0^1(U)\hookrightarrow L^p(U) \quad\text{for }1\le p\le 2^*.
\end{align*}
Thus:
\begin{itemize}
\item for $p\le 2^*$, the embedding holds;
\item for $p>2^*$, generally it fails.
\end{itemize}

---

\item \textbf{5. Dimension matters}

---

\begin{description}
\item If $n>2$,
\begin{align*}
H_0^1(U)\hookrightarrow L^p(U)
\quad\text{for }1\le p\le \frac{2n}{n-2}.
\end{align*}

\item If $n=2$,
\begin{align*}
H_0^1(U)\hookrightarrow L^p(U)
\quad\forall p<\infty.
\end{align*}
No finite critical exponent exists.

\item If $n=1$,
\begin{align*}
H_0^1(U)\hookrightarrow C^0(\overline U)
\end{align*}
and in particular
\begin{align*}
H_0^1(U)\hookrightarrow L^\infty(U).
\end{align*}
So in low dimensions, Sobolev functions are even continuous.
\end{description}

\item \textbf{6. Why this matters in PDE}

Suppose
\begin{align*}
-\Delta u=f(u).
\end{align*}
To make sense of the weak formulation, we need $f(u)$ to be integrable.

If $u\in H_0^1(U)$, Sobolev embedding gives:
\begin{align*}
u\in L^p(U),
\end{align*}
which gives integrability for nonlinear terms like:
\begin{align*}
u^q,\quad |u|^{q-1}u.
\end{align*}
This is essential in nonlinear Partial Differential Equations.

\item \textbf{7. Compact Sobolev embedding}

There is also the \DEFINE{Rellich–Kondrachov theorem}:

If
\begin{align*}
1\le p<2^*,
\end{align*}
then
\begin{align*}
H_0^1(U)\hookrightarrow L^p(U)
\end{align*}
is \DEFINE{compact}.

That means:

If $u_j$ is bounded in $H_0^1(U)$, then some subsequence converges strongly in $L^p(U)$.

This compactness is one of the main tools for proving existence of weak solutions.

\item \textbf{8. Intuition}

If $\nabla u\in L^2$, then $u$ cannot oscillate too wildly.

That extra regularity forces better behavior of $u$, such as:
\begin{itemize}
\item better integrability,
\item continuity,
\item compactness.
\end{itemize}
So:

\> \DEFINE{Derivative bounds control the function itself.}

That is the heart of Sobolev embedding.

\item \textbf{Final summary}

The Sobolev Embedding Theorem says:

\> If a function has sufficiently integrable derivatives, then the function belongs to stronger spaces.

For example, for $n>2$,
\begin{align*}
H_0^1(U)\hookrightarrow L^p(U) \quad\text{for }1\le p\le \frac{2n}{n-2}.
\end{align*}
This means:
\begin{align*}
\|u\|_{L^p(U)} \le C\|u\|_{H_0^1(U)}.
\end{align*}

This theorem is fundamental because it turns derivative estimates into function estimates, which is essential in PDE theory.
\end{description}

If you'd like, I can explain \DEFINE{where the critical exponent $\frac{2n}{n-2}$ comes from}, which is the key scaling argument behind Sobolev embedding.

%%%%%%%%%%%%%%%
\newpage
\item \DEFINE{Explain Calculus of Variations}

The \DEFINE{calculus of variations} is the branch of mathematics where we seek a \DEFINE{function} that minimizes or maximizes a quantity called a \DEFINE{functional}.

Instead of minimizing numbers like
\begin{align*}
f(x),
\end{align*}
we minimize expressions like
\begin{align*}
J[u].
\end{align*}
This is fundamental in Calculus of Variations and is one of the main tools in Partial Differential Equations.

\item \textbf{1. Ordinary calculus vs calculus of variations}

In ordinary calculus:
\begin{align*}
f(x)=x^2+1
\end{align*}
and we look for $x$ minimizing $f(x$).

In calculus of variations:
\begin{align*}
J[u]=\int_a^b F(x,u,u')\,dx
\end{align*}
and we look for a \DEFINE{function $u(x)$} minimizing $J$.

So:
	\begin{itemize}
    \item ordinary calculus: unknown is a number,
    \item calculus of variations: unknown is a function.
	\end{itemize}

\item \textbf{2. Example of a functional}

A typical functional is:
\begin{align*}
J[u]=\int_a^b (u')^2\,dx.
\end{align*}
This measures the ``energy" of the function.

We may want to find $u$ minimizing $J[u]$ among all functions with fixed endpoints:
\begin{align*}
u(a)=A,\qquad u(b)=B.
\end{align*}
The minimizer is the straight line between $A$ and $B$.

\item \textbf{3. The main idea: perturb the function}

Suppose $u$ is a minimizer.

Take a small perturbation:
\begin{align*}
u+\varepsilon v,
\end{align*}
where $v$ is a test function.

Define:
\begin{align*}
\phi(\varepsilon)=J[u+\varepsilon v].
\end{align*}
If $u$ minimizes $J$, then
\begin{align*}
\phi'(0)=0.
\end{align*}
This gives the \DEFINE{Euler–Lagrange equation}, the necessary condition for minimizers.

\item \textbf{4. Euler–Lagrange equation}

For
\begin{align*}
J[u]=\int_a^b F(x,u,u')\,dx,
\end{align*}
the minimizer satisfies
\begin{align*}
\frac{\partial F}{\partial u} - \frac{d}{dx} \left( \frac{\partial F}{\partial u'} \right) =0.
\end{align*}
This is the Euler–Lagrange equation.

It is the analogue of
\begin{align*}
f'(x)=0
\end{align*}
from ordinary calculus.

\item \textbf{5. Example}

Minimize
\begin{align*}
J[u]=\int_a^b (u')^2\,dx.
\end{align*}
Then
\begin{align*}
F=(u')^2.
\end{align*}

So:
\begin{align*}
\frac{\partial F}{\partial u}=0, \qquad \frac{\partial F}{\partial u'}=2u'.
\end{align*}
Thus,
\begin{align*}
\frac{d}{dx}(2u')=0,
\end{align*}
which gives
\begin{align*}
u''=0.
\end{align*}
So the minimizer is linear.

\item \textbf{6. Connection to PDEs}

Many PDEs arise as Euler–Lagrange equations.

For example, consider
\begin{align*}
J[u]=\int_U \frac12 |\nabla u|^2-fu\,dx.
\end{align*}
Its Euler–Lagrange equation is
\begin{align*}
-\Delta u=f.
\end{align*}
So solving the PDE is equivalent to minimizing the energy.

This is why variational methods are so important in Partial Differential Equations.

\item \textbf{7. Direct method in calculus of variations}

A standard way to prove existence of minimizers:

	\begin{description}
    \item Step 1: Choose a minimizing sequence

Take $u_j$ such that
\begin{align*}
J[u_j]\to \inf J.
\end{align*}

    \item Step 2: Show boundedness

Prove $(u_j)$ is bounded in some Sobolev space.

    \item Step 3: Extract a weakly convergent subsequence

Use weak compactness:
\begin{align*}
u_j\rightharpoonup u.
\end{align*}

    \item Step 4: Use lower semicontinuity

Show
\begin{align*}
J[u]\le \liminf J[u_j].
\end{align*}
Then $u$ minimizes $J$.
	\end{description}

This is the \DEFINE{direct method}, and it is exactly what you used in the Sobolev minimization problem.

\item \textbf{8. Physical meaning}

Many physical systems minimize an energy:
	\begin{itemize}
    \item elastic membranes minimize elastic energy,
    \item geodesics minimize length,
    \item soap films minimize area.
	\end{itemize}
The calculus of variations gives the equations governing these minimizers.

Examples:
	\begin{itemize}
    \item geodesics,
    \item minimal surfaces,
    \item Laplace equation,
    \item elasticity,
    \item quantum mechanics.
	\end{itemize}

\item \textbf{9. Why it matters}

The calculus of variations turns PDE problems into minimization problems.

This is useful because minimization problems can often be solved using:
	\begin{itemize}
    \item weak compactness,
    \item Sobolev embeddings,
    \item lower semicontinuity.
	\end{itemize}

Thus:

\> \DEFINE{Variational methods provide existence results even when explicit solutions are impossible.}

\item \textbf{Final summary}

The calculus of variations studies how to minimize functionals like
\begin{align*}
J[u]=\int F(x,u,\nabla u)\,dx.
\end{align*}
The minimizers satisfy Euler–Lagrange equations, often PDEs.

It is a central framework for:
\begin{enumerate}
\item deriving PDEs,
\item proving existence of solutions,
\item interpreting solutions as energy minimizers.
\end{enumerate}

If you'd like, I can show how the minimization problem
\begin{align*}
\inf_{u\in M}\int_U |\nabla u|^2\,dx
\end{align*}
leads to the nonlinear PDE
\begin{align*}
-\Delta u=\lambda |u|^{p-2}u
\end{align*}
via the Euler–Lagrange equation.

\newpage
	\item \DEFINE{Explain the `energy argument'.}


The \DEFINE{energy argument} is a method for proving properties of solutions of PDEs—especially \DEFINE{uniqueness}, \DEFINE{stability}, and \DEFINE{a priori estimates}—by introducing an ``energy" quantity and studying how it changes in time.

It is one of the most important tools in Partial Differential Equations.

\begin{description}

\item \textbf{1. Basic idea}

You define an energy functional $E(t)$ that measures the size of the solution at time $t$.

Typical form:
\begin{align*}
E(t)=\text{kinetic energy}+\text{potential energy}.
\end{align*}
Then:
\begin{enumerate}
\item differentiate $E(t$),
\item use the PDE to simplify $E'(t)$,
\item show $E'(t)\le 0$, or $E'(t)=0$,
\item conclude bounds or uniqueness.
\end{enumerate}

\item \textbf{2. Example: wave equation}

Consider
\begin{align*}
u_{tt}-\Delta u=0.
\end{align*}
Define the energy:
\begin{align*}
E(t)=\frac12\int_U \left(u_t^2+|\nabla u|^2\right)\,dx.
\end{align*}
Interpretation:
\begin{itemize}
    \item $u_t^2$: kinetic energy,
    \item $|\nabla u|^2$: potential energy.
\end{itemize}

\item \textbf{3. Differentiate the energy}

Differentiate:
\begin{align*}
E'(t) = \int_U u_tu_{tt}\,dx+\int_U \nabla u\cdot\nabla u_t\,dx.
\end{align*}
Integrate the second term by parts:
\begin{align*}
\int_U \nabla u\cdot\nabla u_t = -\int_U (\Delta u)u_t +\int_{\partial U}\frac{\partial u}{\partial\nu}u_t.
\end{align*}
So
\begin{align*}
E'(t) = \int_U u_t(u_{tt}-\Delta u)\,dx +\int_{\partial U}\frac{\partial u}{\partial\nu}u_t\,dS.
\end{align*}
Now use the PDE:
\begin{align*}
u_{tt}-\Delta u=0.
\end{align*}
Then
\begin{align*}
E'(t)=\int_{\partial U}\frac{\partial u}{\partial\nu}u_t\,dS.
\end{align*}
So the PDE tells us how energy changes.

\item \textbf{4. Why this helps}

Depending on the boundary condition:
	\begin{itemize}
    \item If $\partial_\nu u=0$,
\begin{align*}
E'(t)=0.
\end{align*}
Energy is conserved.

    \item If $u_t+\partial_\nu u=0$,

then
\begin{align*}
\partial_\nu u=-u_t,
\end{align*}
so
\begin{align*}
E'(t) = -\int_{\partial U}u_t^2\,dS\le 0.
\end{align*}
Energy decreases.

	\end{itemize}
Thus the PDE + boundary conditions control the energy.

\item \textbf{5. Energy arguments prove uniqueness}

Suppose two solutions $u_1,u_2$ exist.

Set
\begin{align*}
w=u_1-u_2.
\end{align*}
Then $w$ satisfies the homogeneous PDE with zero initial data.

Define energy:
\begin{align*}
E_w(t)=\frac12\int_U (w_t^2+|\nabla w|^2)\,dx.
\end{align*}
Then often:
\begin{align*}
E_w'(t)\le 0.
\end{align*}
Since initial data are zero,
\begin{align*}
E_w(0)=0.
\end{align*}
Thus,
\begin{align*}
E_w(t)\equiv 0.
\end{align*}
Hence
\begin{align*}
w_t=0,\quad \nabla w=0,
\end{align*}
so
\begin{align*}
w\equiv 0.
\end{align*}
Therefore the solution is unique.

This is the core of the energy argument.

\item \textbf{6. Energy arguments give bounds}

Sometimes one gets:
\begin{align*}
E'(t)\le CE(t).
\end{align*}
Then Grönwall's inequality gives:
\begin{align*}
E(t)\le E(0)e^{Ct}.
\end{align*}
This yields bounds like
\begin{align*}
|u(t)|\le Ce^{Ct}.
\end{align*}
These are \DEFINE{a priori estimates}.

\item \textbf{7. Why the name ``energy"?}

Because in physical systems the energy often has the form:
\begin{align*}
\text{kinetic}+\text{potential}.
\end{align*}
For example:
\begin{align*}
E(t)=\frac12\int (u_t^2+|\nabla u|^2)\,dx
\end{align*}
is exactly the physical energy of a vibrating string or membrane.

So the PDE describes the evolution of energy.

\item \textbf{8. General structure of an energy argument}

\begin{description}

    \item Step 1:

Define an energy $E(t)$.

    \item Step 2:

Differentiate:
\begin{align*}
E'(t).
\end{align*}

    \item Step 3:

Use the PDE to rewrite $E'(t)$.

    \item Step 4:

Show:
\begin{align*}
E'(t)\le 0
\quad\text{or}\quad
E'(t)\le CE(t).
\end{align*}

    \item Step 5:

Conclude uniqueness or bounds.

\end{description}

\item \textbf{Final summary}

The \DEFINE{energy argument} is:

\> Define an energy measuring the size of the solution, compute its time derivative using the PDE, and use the resulting inequality to prove uniqueness or estimates.

For uniqueness:

\begin{align*}
E(0)=0,\quad E'(t)\le 0
\implies E(t)=0
\implies u=0.
\end{align*}

This is one of the most powerful methods in PDE theory.

\end{description}

If you'd like, I can show how the \DEFINE{energy argument changes for the heat equation}, where the energy decreases because diffusion dissipates energy.


%%%%%%%%%%%%%%%%%%%%%%%%%%%%%
\newpage
	\item \textit{What are the energy formulas for the standard PDEs}.


For the ``standard" PDEs, the \DEFINE{energy formula} is the identity (or inequality) satisfied by a natural norm of the solution, usually obtained by:
\begin{enumerate}
\item multiplying the PDE by the solution (or $u_t)$,
\item integrating over the domain,
\item integrating by parts.
\end{enumerate}

These formulas measure conservation or dissipation of energy and are fundamental in Partial Differential Equations.

\begin{description}
\item \textbf{1. Heat equation}

Consider
\begin{align*}
u_t-\Delta u=0 \quad\text{in }U.
\end{align*}
Multiply by $u$:
\begin{align*}
uu_t-u\Delta u=0.
\end{align*}
Integrate:
\begin{align*}
\int_U uu_t\,dx-\int_U u\Delta u\,dx=0.
\end{align*}
Using integration by parts,
\begin{align*}
\int_U u\Delta u\,dx=-\int_U |\nabla u|^2\,dx
\end{align*}
(for zero boundary data), so:
\begin{align*}
\int_U uu_t\,dx+\int_U |\nabla u|^2\,dx=0.
\end{align*}
Since
\begin{align*}
\frac{d}{dt}\frac12\int_U u^2\,dx = \int_U uu_t\,dx,
\end{align*}
we obtain
\begin{align*}
\boxed{
\frac{d}{dt}\frac12\int_U u^2\,dx + \int_U |\nabla u|^2\,dx = 0.  }
\end{align*}

\textbf{Interpretation:}

\begin{align*}
E(t)=\frac12\int_U u^2\,dx
\end{align*}
is decreasing:
\begin{align*}
E'(t)=-\int_U |\nabla u|^2\,dx\le 0.
\end{align*}
So the heat equation \DEFINE{dissipates energy}.

\item \textbf{2. Wave equation}

Consider
\begin{align*}
u_{tt}-\Delta u=0.
\end{align*}
Multiply by $u_t$:
\begin{align*}
u_tu_{tt}-u_t\Delta u=0.
\end{align*}
Integrate:
\begin{align*}
\int_U u_tu_{tt}\,dx-\int_U u_t\Delta u\,dx=0.
\end{align*}
Integrate by parts:
\begin{align*}
-\int_U u_t\Delta u\,dx=\int_U \nabla u\cdot\nabla u_t\,dx.
\end{align*}
Thus,
\begin{align*}
\int_U u_tu_{tt}\,dx+\int_U \nabla u\cdot\nabla u_t\,dx=0.
\end{align*}
Recognize:
\begin{align*}
\frac{d}{dt}\frac12\int_U u_t^2\,dx = \int_U u_tu_{tt}\,dx,
\end{align*}
and
\begin{align*}
\frac{d}{dt}\frac12\int_U |\nabla u|^2\,dx = \int_U \nabla u\cdot\nabla u_t\,dx.
\end{align*}
Therefore,
\begin{align*}
\boxed{
\frac{d}{dt}\frac12\int_U \left(u_t^2+|\nabla u|^2\right)\,dx=0.
}
\end{align*}

\textbf{Interpretation:}

The energy
\begin{align*}
E(t)=\frac12\int_U(u_t^2+|\nabla u|^2)\,dx
\end{align*}
is conserved.

So the wave equation \DEFINE{preserves energy}.

\item \textbf{3. Poisson equation}

Consider
\begin{align*}
-\Delta u=f.
\end{align*}
Multiply by $u$:
\begin{align*}
-u\Delta u=fu.
\end{align*}
Integrate:
\begin{align*}
-\int_U u\Delta u\,dx=\int_U fu\,dx.
\end{align*}
Integrating by parts gives:
\begin{align*}
\boxed{
\int_U |\nabla u|^2\,dx = \int_U fu\,dx.
}
\end{align*}

This is the energy identity for elliptic equations.

\textbf{Interpretation:}

The Dirichlet energy

\begin{align*}
\int_U |\nabla u|^2\,dx
\end{align*}
is controlled by the forcing term.

\item \textbf{4. Reaction-diffusion equation}

Consider
\begin{align*}
u_t-\Delta u+cu=0.
\end{align*}
Multiply by $u$:
\begin{align*}
uu_t-u\Delta u+cu^2=0.
\end{align*}
Then:
\begin{align*}
\boxed{
\frac{d}{dt}\frac12\int_U u^2\,dx + \int_U |\nabla u|^2\,dx + c\int_U u^2\,dx = 0.
}
\end{align*}
Energy decays even faster because of the reaction term.

\item \textbf{5. Damped wave equation}

Consider
\begin{align*}
u_{tt}-\Delta u+u_t=0.
\end{align*}
Multiply by $u_t$:
\begin{align*}
u_tu_{tt}-u_t\Delta u+u_t^2=0.
\end{align*}
Then:
\begin{align*}
\boxed{
\frac{d}{dt}\frac12\int_U (u_t^2+|\nabla u|^2)\,dx + \int_U u_t^2\,dx = 0.
}
\end{align*}
The energy decreases due to damping.

\item \textbf{6. Schrödinger equation}

Consider
\begin{align*}
iu_t+\Delta u=0.
\end{align*}
\begin{align*}
\boxed{
\frac{d}{dt}\int_U |u|^2\,dx=0.
}
\end{align*}

So the $L^2$-norm is conserved.

\item \textbf{7. General pattern}

For PDEs:

	\begin{itemize}

    \item Diffusion-type equations:

Energy decreases.

\begin{align*}
E'(t)\le 0.
\end{align*}

Examples:

	\begin{itemize}
    \item heat equation,
    \item reaction-diffusion,
    \item damped wave.
	\end{itemize}

    \item Conservative equations:

Energy is constant.

\begin{align*}
E'(t)=0.
\end{align*}

Examples:

	\begin{itemize}
    \item wave equation,
    \item Schrödinger equation.
	\end{itemize}

    \item Elliptic equations:

Energy gives static bounds.

\begin{align*}
\int |\nabla u|^2 = \int fu.
\end{align*}
	\end{itemize}

---

\item \textbf{Final summary}

The standard energy formulas are:

	\begin{itemize}
    \item Heat equation:

\begin{align*}
\boxed{
\frac{d}{dt}\frac12\int u^2 + \int |\nabla u|^2 = 0
}
\end{align*}
Energy decays.

    \item Wave equation:

\begin{align*}
\boxed{
\frac{d}{dt}\frac12\int (u_t^2+|\nabla u|^2)=0
}
\end{align*}
Energy conserved.

    \item Poisson equation:

\begin{align*}
\boxed{
\int |\nabla u|^2=\int fu
}
\end{align*}
Static energy balance.

    \item Damped wave:

\begin{align*}
\boxed{
\frac{d}{dt}E(t)+\int u_t^2=0
}
\end{align*}
Energy decays.

These are the main ``energy identities" used to prove:
\begin{enumerate}
\item uniqueness,
\item stability,
\item a priori bounds,
\item existence.
\end{enumerate}
	\end{itemize}

\end{description}

If you'd like, I can derive the \DEFINE{energy formula for the general second-order hyperbolic PDE}
\begin{align*}
u_{tt}-\operatorname{div}(A\nabla u)+cu=0,
\end{align*}
which unifies the wave and damped wave cases.


\end{description} %%% end of Questions

\end{document}
